% Clase 11 --- El condensador con dieléctrico (§4).
% (a) De dónde sale Q': la gaussiana de siempre encierra la carga de la placa Y
%     la carga de polarización de la cara vecina del aislante. La ley de Gauss
%     usual exige TODA la carga encerrada, y de ahí sale E = (Q+Q')/(eps0 A).
% (b) Meter el dieléctrico no es una sola operación: según se deje fija DeltaV
%     (conectado a la fuente) o Q (desconectado), cambia una cosa u otra.
\begin{center}
\begin{tikzpicture}[scale=1]

  % =============== (a) la gaussiana ve las dos cargas ===============
  \begin{scope}
    % placas
    \draw[figblue, line width=1.6pt] (-2.0,1.9) -- (2.0,1.9);
    \foreach \x in {-1.7,-1.1,...,1.8} {\node[figlbl,figblue,scale=0.7] at (\x,2.1) {$-$};}
    \node[figlbl, figblue, anchor=west] at (2.15,2.0) {$-Q$};
    \draw[figred, line width=1.6pt] (-2.0,-1.9) -- (2.0,-1.9);
    \foreach \x in {-1.7,-1.1,...,1.8} {\node[figlbl,figred,scale=0.7] at (\x,-2.1) {$+$};}
    \node[figlbl, figred, anchor=west] at (2.15,-2.0) {$+Q$};

    % dieléctrico y sus cargas de polarización
    \fill[figgray!8] (-1.85,-1.5) rectangle (1.85,1.5);
    \draw[figgray, line width=0.7pt] (-1.85,-1.5) rectangle (1.85,1.5);
    \foreach \x in {-1.5,-0.9,...,1.6} {\node[figlbl,figred,scale=0.7] at (\x,1.35) {$+$};}
    \foreach \x in {-1.5,-0.9,...,1.6} {\node[figlbl,figblue,scale=0.7] at (\x,-1.35) {$-$};}
    \node[figlblaux, anchor=west] at (1.95,0.6) {$K_E$};

    % gaussiana (cilindro chato) que abraza la placa + y la cara vecina
    \draw[figgray, dashed, line width=0.9pt] (-1.0,-2.25) rectangle (1.0,-1.05);
    \node[figlblaux, anchor=north, align=center] at (0,-2.35)
      {gaussiana: encierra $+Q$ \emph{y} $Q'$};
    \draw[figgray, line width=0.4pt] (-1.55,-1.35) -- (-2.05,-1.35);
    \node[figlbl, figblue, anchor=east] at (-2.1,-1.35) {$Q' < 0$};

    % campo
    \draw[figvecres] (0,-0.9) -- (0,0.9);
    \node[figlbl, figred, anchor=west] at (0.1,0.35) {$\vec E$};

    \node[figlbl, anchor=north, align=center] at (0,-3.15)
      {(a) la ley de Gauss de siempre\\[-0.2em]
       \footnotesize pide \emph{toda} la carga encerrada};
  \end{scope}

  % =============== (b) las dos maneras de meterlo ===============
  \begin{scope}[xshift=8.0cm]
    % --- fila superior: DeltaV fija ---
    \begin{scope}[yshift=1.15cm]
      \draw[figblue, line width=1.3pt] (-1.5,0.55) -- (1.5,0.55);
      \draw[figred, line width=1.3pt] (-1.5,-0.55) -- (1.5,-0.55);
      \fill[figgray!12] (-1.4,-0.42) rectangle (1.4,0.42);
      \node[figlblaux] at (0,0) {$K_E$};
      \node[figlbl, anchor=east, align=right] at (-1.7,0)
        {\footnotesize conectado\\[-0.2em] \footnotesize a la fuente};
      \node[figlbl, anchor=west, align=left] at (1.7,0)
        {$\Delta V$ fija:\\[-0.15em] \textbf{sube} $Q$};
    \end{scope}
    % --- fila inferior: Q fija ---
    \begin{scope}[yshift=-1.15cm]
      \draw[figblue, line width=1.3pt] (-1.5,0.55) -- (1.5,0.55);
      \draw[figred, line width=1.3pt] (-1.5,-0.55) -- (1.5,-0.55);
      \fill[figgray!12] (-1.4,-0.42) rectangle (1.4,0.42);
      \node[figlblaux] at (0,0) {$K_E$};
      % bornes cortados
      \draw[figgray, line width=0.9pt] (-1.5,0.55) -- (-1.9,0.55);
      \draw[figgray, line width=0.9pt] (-1.5,-0.55) -- (-1.9,-0.55);
      \draw[figred, line width=0.9pt] (-2.05,0.4) -- (-1.75,0.7);
      \draw[figred, line width=0.9pt] (-2.05,-0.7) -- (-1.75,-0.4);
      \node[figlbl, anchor=east, align=right] at (-2.2,0)
        {\footnotesize desconectado};
      \node[figlbl, anchor=west, align=left] at (1.7,0)
        {$Q$ fija:\\[-0.15em] \textbf{baja} $\Delta V$};
    \end{scope}

    \node[figlbl, anchor=north, align=center] at (0.2,-3.15)
      {(b) el mismo dieléctrico,\\[-0.2em] dos experimentos\\[0.1em]
       \footnotesize en ambos $C$ crece $\times K_E$};
  \end{scope}

\end{tikzpicture}\\[0.5em]
{\small\itshape El panel (a) es la razón por la que aparece $Q'$ en la cuenta:
la superficie gaussiana no distingue entre carga libre y carga de polarización,
las encierra a las dos, y la ley de Gauss ---que sigue valiendo tal cual con
dieléctricos--- exige usar la suma. El resultado, $Q' = Q(1/K_E - 1)$, tiene el
signo que debía: el paréntesis es negativo porque $K_E > 1$. El panel (b) marca
una distinción que suele confundirse: la capacitancia $C = \varepsilon A/d$
crece en un factor $K_E$ en los dos casos ---es geometría más material, no
depende del experimento---, pero lo que se observa es distinto según qué se
mantenga fijo. Con la fuente conectada entra más carga; desconectado, la carga
no tiene por dónde entrar y lo que cae es el voltaje.}
\end{center}
